% =========================================================
\chapter{Solução Proposta}
% =========================================================

Este capítulo deve responder essencialmente à pergunta:

\begin{center}
	\textbf{O que foi desenvolvido e como funciona?}
\end{center}

% ---------------------------------------------------------
\section{Arquitetura holística do sistema}
% ---------------------------------------------------------

\begin{enumerate}
	\item Qual é a arquitetura geral da solução?
	\item Quais são os principais componentes?
	\item Qual é a função de cada componente?
	\item Como os componentes se comunicam?
	\item Quais são as entradas do sistema?
	\item Quais são os processos internos?
	\item Quais são as saídas?
	\item Quais interfaces existem entre os subsistemas?
	\item Quais componentes são físicos, computacionais, matemáticos ou experimentais?
	\item Como a arquitetura proposta se diferencia das arquiteturas existentes?
\end{enumerate}

{Recomendação.}
Apresentar um diagrama de blocos geral antes de entrar nos detalhes dos
componentes individuais.

% ---------------------------------------------------------
\section{Modelo matemático do sistema}
% ---------------------------------------------------------

\begin{enumerate}
	\item Qual fenômeno ou comportamento está sendo modelado?
	\item Quais são as variáveis de entrada?
	\item Quais são as variáveis de saída?
	\item Quais são os parâmetros do sistema?
	\item Quais leis físicas ou princípios matemáticos governam o sistema?
	\item Quais equações descrevem o comportamento do sistema?
	\item Quais hipóteses foram adotadas?
	\item Quais simplificações foram necessárias?
	\item Quais condições iniciais e de contorno foram consideradas?
	\item O modelo pode ser validado experimentalmente ou comparado com modelos conhecidos?
\end{enumerate}

{Sequência recomendada.}

\begin{center}
	\textbf{
		Fenômeno físico
		$\rightarrow$
		hipóteses
		$\rightarrow$
		variáveis
		$\rightarrow$
		equações
		$\rightarrow$
		parâmetros
		$\rightarrow$
		condições
		$\rightarrow$
		modelo final.
	}
\end{center}

% ---------------------------------------------------------
\section{Fluxograma do sistema}
% ---------------------------------------------------------

\begin{enumerate}
	\item Qual é o estado inicial do sistema?
	\item Qual é o primeiro processamento realizado?
	\item Quais decisões precisam ser tomadas?
	\item Quais condições determinam cada caminho?
	\item Existem ciclos ou realimentações?
	\item Existem processos iterativos?
	\item Quais dados entram em cada etapa?
	\item Quais dados saem de cada etapa?
	\item Qual é o critério de parada?
	\item Qual é o resultado final do algoritmo ou processo?
\end{enumerate}

O fluxograma deve permitir que um terceiro pesquisador reconstrua a lógica
operacional da solução.

% ---------------------------------------------------------
\section{Considerações finais}
% ---------------------------------------------------------

\begin{enumerate}
	\item Qual solução foi efetivamente desenvolvida?
	\item Quais são seus principais componentes?
	\item Quais são suas principais características?
	\item Quais hipóteses foram adotadas?
	\item Quais são suas limitações?
	\item Quais parâmetros serão avaliados?
	\item Como a solução será testada?
	\item Quais critérios serão utilizados para determinar se ela funciona?
	\item Como ela será comparada às soluções da literatura?
	\item O que precisa ser demonstrado experimentalmente ou por simulação?
\end{enumerate}

